\documentclass[../../数学分析培优讲义(答案版).tex]{subfiles}

\begin{document}

\setcounter{chapter}{2}
\pagestyle{fancy}

\chapter{一元函数的连续性}

\section{连续性证明}

\begin{example}
研究函数
\[
    f(x)=\lim\limits_{n\to\infty}\dfrac{x^n-1}{x^n+1}
\]
的连续性.
\end{example}

\begin{proof}[解]
当 $|x|<1$ 时，$x^n\to0$，故极限为 $-1$；当 $|x|>1$ 时，分子分母同除以 $x^n$，极限为 $1$；当 $x=1$ 时，极限为 $0$。在 $x=-1$ 处，奇数项的分母为零，原数列没有定义。因此
\[
f(x)=
\begin{cases}
    -1, & -1<x<1,\\
    0, & x=1,\\
    1, & x<-1\text{ 或 }x>1.
\end{cases}
\]
由此可见，$f$ 在其定义域内除 $x=1$ 外处处连续；在 $x=1$ 处左右极限分别为 $-1$ 与 $1$，故不连续。$x=-1$ 不属于定义域，且该点两侧极限也不相等，因而不能连续延拓到该点。
\end{proof}

\begin{example}[\markstar]
证明 Riemann 函数
\[
R(x)=
\begin{cases}
    \dfrac{1}{q}, & x=\dfrac{p}{q}\text{ 为既约分数},\ q>0,\\
    0, & x\text{ 为无理数},
\end{cases}
\]
在无理点上连续,在有理点上间断.
\end{example}

\begin{proof}[解]
先证无理点处连续。设 $x_0$ 为无理数。给定 $\varepsilon>0$，取正整数 $N$ 使 $N>1/\varepsilon$。任一有界区间内，分母不超过 $N$ 的既约分数只有有限个，故可取 $\delta>0$，使 $(x_0-\delta,x_0+\delta)$ 中不含这些分数。于是当 $|x-x_0|<\delta$ 时，若 $x$ 为无理数，则 $R(x)=0$；若 $x=p/q$ 为既约分数，则 $q>N$，从而
\[
    |R(x)-R(x_0)|=R(x)=\dfrac{1}{q}<\varepsilon.
\]
所以 $R$ 在 $x_0$ 连续。

若 $x_0=p_0/q_0$ 为既约有理数，则 $R(x_0)=1/q_0>0$。取无理数列 $x_n\to x_0$，有 $R(x_n)=0$，故 $R(x_n)\not\to R(x_0)$。因此 $R$ 在每个有理点间断。
\end{proof}

\begin{example}
讨论函数
\[
f(x)=
\begin{cases}
    x(1-x), & x\text{ 为有理数},\\
    x(1+x), & x\text{ 为无理数},
\end{cases}
\]
的连续性.
\end{example}

\begin{proof}[解]
有理数与无理数在实轴上都稠密。若 $f$ 在 $x_0$ 连续，则分别沿有理数列和无理数列趋于 $x_0$，必有
\[
    x_0(1-x_0)=x_0(1+x_0),
\]
即 $x_0=0$。反之，当 $x\to0$ 时，两种情形均有
\[
    |f(x)|\leq |x|(1+|x|)\to0=f(0),
\]
故 $f$ 在 $x=0$ 连续。因此唯一连续点为 $x=0$。
\end{proof}

\begin{example}
设
\[
f(x)=
\begin{cases}
    1, & x\geq0,\\
    -1, & x<0,
\end{cases}
\qquad g(x)=\sin x,
\]
讨论 $f(g(x))$ 的连续性.
\end{example}

\begin{proof}[解]
复合函数为
\[
f(g(x))=
\begin{cases}
    1, & \sin x\geq0,\\
    -1, & \sin x<0.
\end{cases}
\]
在每个不满足 $\sin x=0$ 的点，$\sin x$ 在该点邻域内符号不变，故复合函数局部为常数，从而连续。对于 $x=k\pi$，其左右两侧 $\sin x$ 异号，复合函数的左右极限分别为 $1$ 与 $-1$，故在这些点间断。于是间断点恰为 $k\pi$，其中 $k\in\mathbb{Z}$。
\end{proof}

\begin{example}
设 $f(x)$ 在 $(0,1)$ 内有定义,且函数 $e^x f(x)$ 与 $e^{-f(x)}$ 在 $(0,1)$ 内都是单调不减的,试证:$f(x)$ 在 $(0,1)$ 内连续.
\end{example}

\begin{proof}[解]
由 $e^{-f(x)}$ 单调不减可知 $f$ 单调不增。任取 $x_0\in(0,1)$，记 $f(x_0-)$ 与 $f(x_0+)$ 为其左右极限，则
\[
    f(x_0-)\geq f(x_0)\geq f(x_0+).
\]
另一方面，$e^x f(x)$ 单调不减，故
\[
e^{x_0}f(x_0-)\leq e^{x_0}f(x_0)\leq e^{x_0}f(x_0+).
\]
因 $e^{x_0}>0$，上式给出
\[
    f(x_0-)\leq f(x_0)\leq f(x_0+).
\]
两组不等式合并即得 $f(x_0-)=f(x_0)=f(x_0+)$。由于 $x_0$ 任意，$f$ 在 $(0,1)$ 内连续。
\end{proof}

\begin{example}
设 $f(x)$ 在 $(a,b)$ 上至多只有第一类间断点,且对任意 $x,y\in(a,b)$,
\[
    f\left(\dfrac{x+y}{2}\right)\leq\dfrac{f(x)+f(y)}{2}.
\]
证明:$f(x)$ 在 $(a,b)$ 上连续.
\end{example}

\begin{proof}[解]
题设不等式说明 $f$ 为中点凸函数。由“至多只有第一类间断点”可知，任一点附近 $f$ 都有界，特别地，$f$ 在某个非退化闭区间上有上界。中点凸函数在某个区间上局部有上界时必连续，下面给出所需估计。

固定 $x_0\in(a,b)$，取 $r>0$ 使 $\lbrack x_0-r,x_0+r\rbrack\subset(a,b)$，并设该区间上 $f\leq M$。反复使用中点凸性可得对任意二进有理数 $\lambda\in\lbrack0,1\rbrack$，
\[
    f((1-\lambda)x+\lambda y)
    \leq(1-\lambda)f(x)+\lambda f(y).
\]
当 $0<|h|<r/2$ 时，可取形如 $\lambda=2^{-m}$ 的二进有理数，使
\[
    \dfrac{|h|}{r}\leq\lambda<\dfrac{2|h|}{r}.
\]
令 $u=x_0+h/\lambda$，则 $|u-x_0|\leq r$，且
$x_0+h=(1-\lambda)x_0+\lambda u$。于是
\[
    f(x_0+h)-f(x_0)\leq\lambda(M-f(x_0))\to0.
\]
这里 $h\to0$ 时 $\lambda<2|h|/r\to0$。再由
\[
    f(x_0)\leq\dfrac{f(x_0+h)+f(x_0-h)}{2}
\]
及对 $-h$ 的同一上界，可得
\[
    f(x_0+h)-f(x_0)
    \geq f(x_0)-f(x_0-h).
\]
对 $-h$ 的上估计表明右端的下极限不小于 $0$；结合已得的上极限不大于 $0$，可知
$f(x_0+h)\to f(x_0)$。$x_0$ 任意，故 $f$ 在 $(a,b)$ 上连续。
\end{proof}

\begin{example}[\markstar]
\begin{enumerate}[label=(\arabic*)]
    \item 证明:若函数 $g(x),f(x)$ 连续,则函数
    \[
        \varphi(x)=\min\{f(x),g(x)\},
        \qquad
        \psi(x)=\max\{f(x),g(x)\}
    \]
    连续;
    \item 设 $f_1(x),f_2(x),f_3(x)$ 在 $\lbrack a,b\rbrack$ 上连续,令函数 $f$ 的值等于三值 $f_1(x),f_2(x),f_3(x)$ 中介于其它两值之间的那个值,证明:$f$ 在 $\lbrack a,b\rbrack$ 上连续.
\end{enumerate}
\end{example}

\begin{proof}[解]
利用恒等式
\[
\min\{u,v\}=\dfrac{u+v-|u-v|}{2},
\qquad
\max\{u,v\}=\dfrac{u+v+|u-v|}{2},
\]
可知连续函数的最小值函数与最大值函数仍连续，这证明了第 (1) 问。

三数 $u,v,w$ 的中间值可写成
\[
    u+v+w-\max\{u,v,w\}-\min\{u,v,w\}.
\]
有限个连续函数的最大值与最小值均连续，故令 $u=f_1(x)$、$v=f_2(x)$、$w=f_3(x)$，即知第 (2) 问中的 $f$ 连续。
\end{proof}

\begin{example}
是否存在 $\lbrack a,b\rbrack$ 上的连续函数:在有理点上取值为无理数,在无理点上取值为有理数?请作出判断并说明理由.
\end{example}

\begin{proof}[解]
不存在。反设存在这样的连续函数 $f$。对每个有理数 $r$，闭集
\[
    F_r=\{x\in\lbrack a,b\rbrack\mid f(x)=r\}
\]
不含任何有理点，因而不可能含有非空开区间，所以 $F_r$ 是无处稠密集。又因为每个无理点的函数值均为有理数，
\[
    \lbrack a,b\rbrack
    =\bigl(\mathbb{Q}\cap\lbrack a,b\rbrack\bigr)
    \cup\bigcup\limits_{r\in\mathbb{Q}}F_r.
\]
右端是可数个无处稠密集的并：单点集无处稠密，而 $\mathbb{Q}\cap\lbrack a,b\rbrack$ 可数。这与完备度量空间 $\lbrack a,b\rbrack$ 的 Baire 定理矛盾。因此所求连续函数不存在。
\end{proof}

\begin{example}
设 $f(x),g(x)$ 是定义在 $(a,b)$ 上的单调函数,且对 $\lbrack a,b\rbrack$ 中的任一有理数 $r$,均有 $f(r)=g(r)$,则 $f(x)$ 与 $g(x)$ 的连续点相同,在连续点上的函数值也相同.
\end{example}

\begin{proof}[解]
只证单调不减的情形；单调不增同理。任取 $x_0\in(a,b)$，由有理数的稠密性可选有理数列 $r_n\uparrow x_0$ 与 $s_n\downarrow x_0$。单调性给出
\[
    f(x_0-)=\lim\limits_{n\to\infty}f(r_n)
    =\lim\limits_{n\to\infty}g(r_n)=g(x_0-),
\]
以及 $f(x_0+)=g(x_0+)$。因此两函数在 $x_0$ 连续的条件相同，即共同的左右极限相等。若 $x_0$ 是连续点，则当 $x_0$ 为有理数时由题设直接有 $f(x_0)=g(x_0)$；当 $x_0$ 为无理数时，单调性迫使两函数在该点的值都等于共同的左右极限。故二者连续点相同，且在连续点取值相同。
\end{proof}

\begin{example}
设 $f(x)$ 定义在 $(a,b)$ 上,且对 $(a,b)$ 中任一点 $x_0$,均存在极限 $\displaystyle\lim\limits_{x\to x_0}f(x)$.若令
\[
    g(x)=\lim\limits_{t\to x}f(t),
    \qquad x\in(a,b),
\]
则 $g(x)\in C((a,b))$.
\end{example}

\begin{proof}[解]
固定 $x_0\in(a,b)$，记
\[
    L=\lim\limits_{t\to x_0}f(t)=g(x_0).
\]
给定 $\varepsilon>0$，可取 $\delta>0$，使 $0<|t-x_0|<\delta$ 时 $|f(t)-L|<\varepsilon$。若 $0<|x-x_0|<\delta/2$，则可取趋于 $x$ 且仍落在上述去心邻域内的点列 $t_n$。于是
\[
    |g(x)-L|
    =\lim\limits_{n\to\infty}|f(t_n)-L|\leq\varepsilon.
\]
故 $g(x)\to g(x_0)$，从而 $g$ 在 $x_0$ 连续。由 $x_0$ 的任意性，$g\in C((a,b))$。
\end{proof}

\subsection{确界函数}

\begin{example}
设 $f(x)$ 在 $\lbrack a,b\rbrack$ 上有界,函数
\[
    M(x)=\sup\limits_{a\leq t\leq x}f(t),
    \qquad
    m(x)=\inf\limits_{a\leq t\leq x}f(t).
\]
则:
\begin{enumerate}[label=(\arabic*)]
    \item $M(x),m(x)$ 在 $\lbrack a,b\rbrack$ 上左连续,并举例说明它们可以不右连续;
    \item 若 $f(x)$ 在 $\lbrack a,b\rbrack$ 上连续,则 $M(x),m(x)$ 在 $\lbrack a,b\rbrack$ 上连续.
\end{enumerate}
\end{example}

\begin{proof}[解]
第 (1) 问按当前题面并不成立。取 $c\in(a,b)$，令
\[
f(t)=
\begin{cases}
    1, & t=c,\\
    0, & t\ne c.
\end{cases}
\]
则 $M(x)=0$（$x<c$）而 $M(x)=1$（$x\geq c$），故 $M$ 在 $c$ 处不左连续。类似地取 $f(c)=0$、$f(t)=1$（$c<t<c+\eta$），可使 $M$ 在 $c$ 处不右连续。因此“任意有界 $f$ 均使 $M,m$ 左连续”这一结论需要核对原始条件。

第 (2) 问成立。若 $f$ 连续，则它在紧区间上一致连续。给定 $\varepsilon>0$，取 $\delta>0$，使 $|u-v|<\delta$ 时 $|f(u)-f(v)|<\varepsilon$。当 $x<y$ 且 $y-x<\delta$ 时，若 $M(y)>M(x)$，则新增的最大值只能在 $\lbrack x,y\rbrack$ 上取得，于是
\[
    0\leq M(y)-M(x)\leq
    \max\limits_{t\in\lbrack x,y\rbrack}|f(t)-f(x)|<\varepsilon.
\]
因 $M$ 单调不减，这已证明 $M$ 连续；对 $m$ 同理。
\end{proof}

\begin{example}
设 $f(x)$ 在 $\lbrack 0,1\rbrack$ 上连续,且 $f(x)>0$,
\[
    R(x)=\sup\limits_{0\leq y\leq x}f(y)\quad(0\leq x\leq1),
    \qquad
    G(x)=\lim\limits_{n\to\infty}\left[\dfrac{f(x)}{R(x)}\right]^n.
\]
试证:当且仅当 $f(x)$ 在 $\lbrack 0,1\rbrack$ 上单增时,$G$ 是连续的.
\end{example}

\begin{proof}[解]
因为 $0<f(x)\leq R(x)$，逐点极限为
\[
G(x)=
\begin{cases}
    1, & f(x)=R(x),\\
    0, & f(x)<R(x).
\end{cases}
\]
若 $f$ 单调不减，则 $R(x)=f(x)$，故 $G\equiv1$，当然连续。

反之，若 $G$ 连续，则 $G$ 是从连通集 $\lbrack0,1\rbrack$ 到离散集 $\{0,1\}$ 的连续函数，因而必为常数。又 $R(0)=f(0)$，所以 $G(0)=1$，从而 $G\equiv1$。于是对每个 $x$ 都有 $f(x)=R(x)$。若 $0\leq x_1<x_2\leq1$，则
\[
    f(x_1)\leq R(x_2)=f(x_2),
\]
故 $f$ 单调不减。
\end{proof}

\begin{example}
已知
\[
f(x)=
\begin{cases}
    x, & 0\leq x<1,\\
    k+1, & k\leq x<k+1,
\end{cases}
\qquad k=1,2,3,\ldots.
\]
求函数
\[
    g(y)=\sup\limits_{f(x)\leq y}x
\]
在 $y\geq0$ 时的具体表达式,并指出 $g(y)$ 在各点处的左右连续性.
\end{example}

\begin{proof}[解]
当 $0\leq y<1$ 时，条件 $f(x)\leq y$ 等价于 $0\leq x\leq y$，故 $g(y)=y$。当 $y\geq1$ 时，设 $n=\lfloor y\rfloor$。所有满足 $k+1\leq y$ 的阶梯区间均被纳入，而区间 $\lbrack n,n+1)$ 尚未被纳入，因此可取的 $x$ 的上确界为 $n$。所以
\[
g(y)=
\begin{cases}
    y, & 0\leq y<1,\\
    \lfloor y\rfloor, & y\geq1.
\end{cases}
\]
$g$ 在 $y=0$ 处右连续，在 $y=1$ 处连续；在每个整数 $n\geq2$ 处右连续而不左连续；在其余点连续。
\end{proof}

\subsection{单射、介值性与单调性的关系}

\subsection{与点集结合 *}

\begin{example}
证明定理:$f(x)$ 在实轴 $x$ 上连续 $\Longleftrightarrow$ 任何开集的逆像仍为开集.
\end{example}

\begin{proof}[解]
设 $f$ 连续，$U$ 为开集。若 $x_0\in f^{-1}(U)$，由于 $U$ 开，可取 $\varepsilon>0$ 使 $(f(x_0)-\varepsilon,f(x_0)+\varepsilon)\subset U$。由连续性，存在 $\delta>0$，使 $|x-x_0|<\delta$ 时 $|f(x)-f(x_0)|<\varepsilon$，故该邻域包含于 $f^{-1}(U)$。所以 $f^{-1}(U)$ 为开集。

反之，若任一开集的逆像均开，则对任意 $x_0$ 与 $\varepsilon>0$，开集 $(f(x_0)-\varepsilon,f(x_0)+\varepsilon)$ 的逆像是包含 $x_0$ 的开集，因而包含某个 $(x_0-\delta,x_0+\delta)$。这正是 $f$ 在 $x_0$ 连续的定义。
\end{proof}

\begin{example}
设 $f(x)$ 在 $(-\infty,+\infty)$ 上有定义,证明:$f(x)$ 连续 $\Longleftrightarrow$ 对任意 $c\in(-\infty,+\infty)$,集合
\[
    \{x\mid f(x)>c\}
    \quad\text{与}\quad
    \{x\mid f(x)<c\}
\]
为开集.
\end{example}

\begin{proof}[解]
若 $f$ 连续，则 $(c,+\infty)$ 与 $(-\infty,c)$ 均为开集，由上一题知其逆像
\[
    \{x\mid f(x)>c\},\qquad \{x\mid f(x)<c\}
\]
均开。

反之，任意开区间 $(\alpha,\beta)$ 的逆像可写为
\[
f^{-1}((\alpha,\beta))
=\{x\mid f(x)>\alpha\}\cap\{x\mid f(x)<\beta\},
\]
故为开集。任意开集都是可数个互不相交开区间的并，其逆像仍开。由上一题的充要条件，$f$ 连续。
\end{proof}

\begin{example}
设 $f:\lbrack 0,1\rbrack\to\mathbb{R}$ 单调递增,且 $f(\lbrack 0,1\rbrack)$ 是闭集,证明:$f$ 在 $\lbrack 0,1\rbrack$ 上连续.
\end{example}

\begin{proof}[解]
本题按当前题面不成立。定义
\[
f(x)=
\begin{cases}
    0, & 0\leq x<\dfrac{1}{2},\\
    1, & \dfrac{1}{2}\leq x\leq1.
\end{cases}
\]
则 $f$ 单调不减，且
\[
    f(\lbrack0,1\rbrack)=\{0,1\}
\]
是闭集，但 $f$ 在 $x=1/2$ 处不连续。因此仅有“值域为闭集”不足以推出连续，题目还需要补充例如值域为区间或介值性等条件。
\end{proof}

\begin{example}
设 $y=f(x)$ 是从 $\lbrack 0,1\rbrack$ 到 $\lbrack 0,1\rbrack$ 的函数,其图像
\[
    E=\{(x,y)\mid y=f(x),\ x\in\lbrack 0,1\rbrack\}
\]
是 $\lbrack 0,1\rbrack\times\lbrack 0,1\rbrack$ 上的闭集,求证:$f$ 是连续函数.
\end{example}

\begin{proof}[解]
任取数列 $x_n\to x_0$。由于 $f(x_n)\in\lbrack0,1\rbrack$，任一子列都含有收敛子列 $f(x_{n_k})\to y$。于是
\[
    (x_{n_k},f(x_{n_k}))\to(x_0,y).
\]
图像 $E$ 为闭集，故 $(x_0,y)\in E$，从而 $y=f(x_0)$。这说明 $\{f(x_n)\}$ 的每个收敛子列都只能收敛到 $f(x_0)$。由于它位于紧集 $\lbrack0,1\rbrack$ 中，遂有整个数列 $f(x_n)\to f(x_0)$。由序列判别法，$f$ 在 $x_0$ 连续；$x_0$ 任意，故 $f$ 连续。
\end{proof}

\begin{example}
设 $y=f(x)$ 为 $X\to Y$ 的连续函数,$F$ 为 $Y$ 轴上的闭集,试证:$f^{-1}(F)$ 为 $X$ 轴上的闭集.
\end{example}

\begin{proof}[解]
闭集的补集 $Y\setminus F$ 为开集。由 $f$ 连续，开集的逆像 $f^{-1}(Y\setminus F)$ 为开集。又
\[
    f^{-1}(Y\setminus F)=X\setminus f^{-1}(F),
\]
故 $f^{-1}(F)$ 的补集为开集，即 $f^{-1}(F)$ 为闭集。
\end{proof}

\newpage

\section{连续性应用}

\subsection{\texorpdfstring{$f(f(x))$}{f(f(x))} 迭代问题}

\begin{example}
是否存在 $\mathbb{R}$ 上连续函数 $f(x)$,使得
\[
    f(f(x))=e^{-x}.
\]
\end{example}

\begin{proof}[解]
不存在。若存在，则 $f(f(x))=e^{-x}$ 为一一映射，因而 $f$ 本身为一一映射。实轴上的连续单射必严格单调。若 $f$ 严格递增，则 $f\circ f$ 严格递增；若 $f$ 严格递减，则两个严格递减函数的复合仍严格递增。无论哪种情形，$f\circ f$ 都严格递增，但 $e^{-x}$ 严格递减，矛盾。
\end{proof}

\begin{example}
设 $(-\infty,+\infty)$ 上的函数满足
\[
    f(f(x))=-x,
\]
则 $f(x)$ 不是连续函数.
\end{example}

\begin{proof}[解]
由 $f(f(x))=-x$ 可知 $f\circ f$ 为一一映射，故 $f$ 为一一映射。若 $f$ 连续，则它在实轴上必严格单调。与上一题相同，无论 $f$ 严格递增还是严格递减，$f\circ f$ 都严格递增；然而 $-x$ 严格递减，矛盾。因此 $f$ 不可能连续。
\end{proof}

\begin{example}
设 $f:\lbrack 0,1\rbrack\to\lbrack 0,1\rbrack$ 为连续函数,且 $f(0)=0,f(1)=1,f(f(x))\equiv x$,证明:$f(x)\equiv x$.
\end{example}

\begin{proof}[解]
由 $f(f(x))=x$ 知 $f$ 为单射。连续单射 $f:\lbrack0,1\rbrack\to\lbrack0,1\rbrack$ 必严格单调；又 $f(0)=0<f(1)=1$，故 $f$ 严格递增。若某个 $x$ 满足 $f(x)>x$，则
\[
    x<f(x)<f(f(x))=x,
\]
矛盾；若 $f(x)<x$，同理有 $x>f(x)>f(f(x))=x$，也矛盾。因此对一切 $x\in\lbrack0,1\rbrack$，均有 $f(x)=x$。
\end{proof}

\begin{example}
设 $f$ 是从 $\mathbb{R}$ 到 $\mathbb{R}$ 的一对一连续映射,有不动点,且满足
\[
    f(2x-f(x))\equiv x,
    \qquad \forall x\in\mathbb{R}.
\]
试证:$f(x)\equiv x\ (\forall x\in\mathbb{R})$.
\end{example}

\begin{proof}[解]
题设等式说明 $f$ 为满射，而题设已给出 $f$ 为单射，故 $f$ 是连续双射并严格单调。若 $f$ 严格递减，则 $x\mapsto2x-f(x)$ 严格递增，左端 $f(2x-f(x))$ 严格递减，不可能等于严格递增的 $x$。故 $f$ 严格递增。

由单射性，等式可写成
\[
    f^{-1}(x)=2x-f(x).
\]
令 $d=f(x)-x$。代入 $x,f(x),f^2(x),\ldots$，归纳可得
\[
    f^n(x)=x+nd,
    \qquad n\in\mathbb{Z}.
\]
设不动点为 $a$。因 $f$ 严格递增，$x>a$ 时所有 $f^n(x)$（$n\in\mathbb{Z}$）都应大于 $a$，而当 $d\ne0$ 时数列 $x+nd$ 随 $n\in\mathbb{Z}$ 取遍两端无界，必有一项越过 $a$，矛盾。$x<a$ 时同理。因此 $d=0$，即 $f(x)=x$。
\end{proof}

\begin{example}
设 $f\in C(( -\infty,+\infty))$ 且
\[
    f(f(f(x)))=x,
    \qquad x\in(-\infty,+\infty),
\]
则 $f(x)\equiv x$.
\end{example}

\begin{proof}[解]
由 $f^3(x)=x$ 可知 $f$ 为单射，连续单射必严格单调。若 $f$ 严格递减，则 $f^3$ 仍严格递减，不可能等于恒等映射，故 $f$ 严格递增。若某点满足 $f(x)>x$，则
\[
    x<f(x)<f^2(x)<f^3(x)=x,
\]
矛盾；$f(x)<x$ 时同理。因此 $f(x)=x$ 对所有 $x\in\mathbb{R}$ 成立。
\end{proof}

\begin{example}
设 $f:\lbrack 0,1\rbrack\to\lbrack 0,1\rbrack$ 连续,且 $f(0)=0,f(1)=1$,定义
\[
    f^n(x)\overset{\triangle}{=}(f\circ f\circ\cdots\circ f)(x)
    \quad(n\text{ 次复合}).
\]
若存在 $n_0$,使得 $f^{n_0}(x)=x\ (0\leq x\leq1)$,则 $f(x)=x$.
\end{example}

\begin{proof}[解]
由 $f^{n_0}=\operatorname{id}$ 知 $f$ 为单射，故连续函数 $f$ 严格单调。端点条件排除了严格递减，所以 $f$ 严格递增。若 $f(x)>x$，反复利用单调性得到
\[
    x<f(x)<f^2(x)<\cdots<f^{n_0}(x)=x,
\]
矛盾；若 $f(x)<x$，同样得到严格递减的闭链，也矛盾。因此 $f(x)=x$。
\end{proof}

\begin{example}
设 $f:\mathbb{R}\to\mathbb{R}$ 连续,且 $f(x)$ 无不动点,
\[
    f^n(x)\overset{\triangle}{=}(f\circ f\circ\cdots\circ f)(x),
\]
则对任意 $x_0\in\mathbb{R}$,数列 $x_n=f^n(x_0)$ 无界.
\end{example}

\begin{proof}[解]
连续函数 $f(x)-x$ 没有零点，所以它在连通集 $\mathbb{R}$ 上恒正或恒负。若 $f(x)>x$，则
\[
    x_0<x_1<x_2<\cdots.
\]
若此数列有上界，则它收敛于某个 $L$；由连续性，
\[
    L=\lim\limits_{n\to\infty}x_{n+1}
    =\lim\limits_{n\to\infty}f(x_n)=f(L),
\]
与无不动点矛盾。因此 $\{x_n\}$ 无上界。若 $f(x)<x$，则 $\{x_n\}$ 严格递减；若它有下界，同样会收敛到不动点，故无下界。两种情形均说明数列无界。
\end{proof}

\begin{example}
设 $f(x)$ 是 $\mathbb{R}$ 上连续函数,且
\[
    \lim\limits_{x\to\infty}f(f(x))=\infty,
\]
证明:
\[
    \lim\limits_{x\to\infty}f(x)=\infty.
\]
\end{example}

\begin{proof}[解]
本题按当前题面不成立。取连续函数 $f(x)=-x$，则
\[
    f(f(x))=x,
\]
故 $\displaystyle\lim\limits_{x\to+\infty}f(f(x))=+\infty$；但
\[
    \lim\limits_{x\to+\infty}f(x)=-\infty,
\]
与所要求的结论相反。因此题目还需补充排除 $f(x)\to-\infty$ 的条件，例如某种单调方向或值域下界条件。
\end{proof}

\subsection{周期函数}

\begin{example}
证明命题 3.2.2 及推论 3.2.2.
\end{example}

\begin{proof}[解]
记正周期集合为 $P$，令 $T_0=\inf P$。先证 $T_0>0$。若 $T_0=0$，可取周期 $T_n\downarrow0$。对任意 $x,y$，取整数 $k_n$ 使 $|y-x-k_nT_n|\leq T_n$，则由周期性与连续性
\[
    f(x)=f(x+k_nT_n)\longrightarrow f(y),
\]
故 $f$ 为常数，矛盾。因此 $T_0>0$。

再取 $T_n\in P$ 且 $T_n\to T_0$。由连续性，
\[
    f(x+T_0)=\lim\limits_{n\to\infty}f(x+T_n)=f(x),
\]
故 $T_0$ 本身是周期，即为最小正周期。

若 $T$ 是任一正周期，作 Euclid 分解 $T=nT_0+r$，其中 $0\leq r<T_0$。周期之差仍为周期，所以若 $r>0$，则 $r$ 是小于 $T_0$ 的正周期，矛盾；故 $r=0$，即 $T=nT_0$。若 $T$ 与 $\alpha T$ 都是正周期，则存在正整数 $m,n$ 使 $T=mT_0$、$\alpha T=nT_0$，从而 $\alpha=n/m\in\mathbb{Q}$。命题与推论均得证。
\end{proof}

\begin{example}
试证明下列命题:
\begin{enumerate}[label=(\arabic*)]
    \item 设 $f\in C(( -\infty,+\infty))$,若 $T_1,T_2$ 皆为 $f$ 的周期,且 $T_1$ 与 $T_2$ 不可公约(即 $\dfrac{T_1}{T_2}\notin\mathbb{Q}$),则 $f(x)$ 为常数;
    \item 设 $f\in C(( -\infty,+\infty))$ 是周期函数,设其正周期全体形成的数集的下确界为 $T_0$,若 $T_0=0$,则 $f(x)\equiv C$(常数).
\end{enumerate}
\end{example}

\begin{proof}[解]
第 (1) 问中，任意整数 $m,n$ 使 $mT_1+nT_2$ 都是 $f$ 的周期。由于 $T_1/T_2\notin\mathbb{Q}$，集合
\[
    \{mT_1+nT_2\mid m,n\in\mathbb{Z}\}
\]
在 $\mathbb{R}$ 中稠密。给定 $x,y$，可取其中的周期 $S_k\to y-x$，于是
\[
    f(y)=\lim\limits_{k\to\infty}f(x+S_k)=f(x).
\]
故 $f$ 为常数。

第 (2) 问中，由下确界为零可取正周期 $T_n\to0$。与上题相同，对任意 $x,y$，选整数 $k_n$ 使 $x+k_nT_n\to y$，利用周期性与连续性得 $f(x)=f(y)$，故 $f$ 为常数。
\end{proof}

\begin{example}
设 $f(x),g(x)$ 各有正周期 $T_1,T_2$,若 $\dfrac{T_1}{T_2}\notin\mathbb{Q}$,则 $h(x)=f(x)+g(x)$ 非周期函数.
\end{example}

\begin{proof}[解]
本题按当前题面不成立，因为没有要求 $f,g$ 非常数，也没有要求 $T_1,T_2$ 是最小正周期。取
\[
    f(x)\equiv0,
    \qquad
    g(x)\equiv0,
\]
并任取不可公约的正数 $T_1,T_2$，例如 $T_1=1$、$T_2=\sqrt2$。此时 $T_1,T_2$ 分别都是 $f,g$ 的周期，但 $h(x)=f(x)+g(x)\equiv0$ 仍是周期函数。因此需补充条件后才能得到预期结论。
\end{proof}

\subsection{Cauchy 方程}

\begin{example}
设 $f\in C(( -\infty,+\infty))$,且
\[
    f(x+y)=f(x)+f(y),
    \qquad \forall x,y\in\mathbb{R},
\]
证明:这个函数方程除零解外只有解 $f(x)=ax$,其中 $a=f(1)>0$.
\end{example}

\begin{proof}[解]
令 $a=f(1)$。由方程先得 $f(0)=0$、$f(-x)=-f(x)$。对整数 $n$ 与有理数 $m/n$，有
\[
    f(n x)=n f(x),
    \qquad
    f\left(\dfrac{m}{n}\right)=\dfrac{m}{n}a.
\]
对任意 $x\in\mathbb{R}$，取有理数列 $r_k\to x$。由连续性，
\[
    f(x)=\lim\limits_{k\to\infty}f(r_k)
    =\lim\limits_{k\to\infty}a r_k=ax.
\]
反之，将 $f(x)=ax$ 代入可得
\[
    f(x+y)=a(x+y)=ax+ay=f(x)+f(y),
\]
故它确实满足方程。

需要指出，题面中的 $a=f(1)>0$ 不能由现有条件推出；$a<0$ 时同样给出非零连续解，例如 $f(x)=-x$。因此非零解应写为 $f(x)=ax$，其中 $a\ne0$。
\end{proof}

\begin{example}
设 $f\in C(( -\infty,+\infty))$,且
\[
    f\left(\dfrac{x+y}{2}\right)=\dfrac{1}{2}[f(x)+f(y)],
    \qquad \forall x,y\in\mathbb{R},
\]
证明:
\[
    f(x)=[f(1)-f(0)]x+f(0).
\]
\end{example}

\begin{proof}[解]
令 $g(x)=f(x)-f(0)$，则
\[
    g\left(\dfrac{x+y}{2}\right)=\dfrac{g(x)+g(y)}{2},
    \qquad g(0)=0.
\]
取 $y=0$、将 $x$ 换成 $2x$，得 $g(2x)=2g(x)$。因此
\[
    g(x+y)=g\left(\dfrac{2x+2y}{2}\right)
    =\dfrac{g(2x)+g(2y)}{2}=g(x)+g(y).
\]
$g$ 连续，由 Cauchy 方程的连续解结论，$g(x)=g(1)x$。于是
\[
    f(x)=[f(1)-f(0)]x+f(0).
\]
\end{proof}

\begin{example}
证明:满足方程
\[
    f(x+y)=f(x)+f(y),
    \qquad \forall x,y>0,
\]
的唯一非零解为 $f(x)=b\log_a x\ (a>0,a,b\text{ 为常数})$.
\end{example}

\begin{proof}[解]
本题按当前题面不成立。方程是关于加法的 Cauchy 方程，而对数函数把乘法化为加法。最简单的反例是
\[
    f(x)=x,
\]
它对任意 $x,y>0$ 满足 $f(x+y)=f(x)+f(y)$，但不可能在整个 $(0,+\infty)$ 上写成 $b\log_a x$。反过来，$b\log_a x$ 一般满足的是
\[
    f(xy)=f(x)+f(y),
\]
而非题面中的加法方程。因此原题很可能应将 $x+y$ 核对为 $xy$；若保留 $x+y$，还需增加连续性等正则条件，连续解应为 $f(x)=cx$。
\end{proof}

\begin{example}
函数方程
\[
    f(x+y)+f(x-y)=2f(x)f(y),
    \qquad \forall x,y\in\mathbb{R},
\]
在实轴 $\mathbb{R}$ 上的非零连续解为 $f(x)=\cos ax$ 或 $f(x)=\operatorname{ch}ax$($a$ 为常数).
\end{example}

\begin{proof}[解]
令 $y=0$。非零解必满足 $f(0)=1$；再令 $x=0$ 可得 $f(-y)=f(y)$，故 $f$ 为偶函数。

连续性事实上会自动提高为二次可微。取充分小的 $a>0$，使
\[
    A=\int_0^a f(t)\,\mathrm{d}t\ne0.
\]
将方程关于 $y$ 在 $\lbrack0,a\rbrack$ 上积分，得到
\[
    2A f(x)=\int_{x-a}^{x+a}f(t)\,\mathrm{d}t.
\]
右端关于 $x$ 可微，故 $f$ 可微；再对原方程逐次求导可知 $f$ 二次可微。对原方程关于 $y$ 求两次导数并令 $y=0$，得
\[
    f''(x)=f''(0)f(x),
    \qquad
    f(0)=1,
    \qquad
    f'(0)=0.
\]
若 $f''(0)=-a^2<0$，解为 $f(x)=\cos ax$；若 $f''(0)=a^2>0$，解为 $f(x)=\operatorname{ch}ax$；若 $f''(0)=0$，则 $f\equiv1$，它也包含在 $a=0$ 的两类写法中。直接代入可验证这些函数均满足原方程。
\end{proof}

\newpage

\section{闭区间上连续函数性质}

\subsection{有界性、最值性}

\begin{example}
设 $f\in C(\lbrack a,b\rbrack)$,若对任意 $x\in\lbrack a,b\rbrack$,存在 $x'\in\lbrack a,b\rbrack$,使得
\[
    |f(x')|\leq\dfrac{|f(x)|}{2},
\]
则存在 $x_0\in\lbrack a,b\rbrack$,使得 $f(x_0)=0$.
\end{example}

\begin{proof}[解]
连续函数 $|f|$ 在紧区间 $\lbrack a,b\rbrack$ 上取得最小值。设
\[
    m=|f(x_*)|=\min\limits_{x\in\lbrack a,b\rbrack}|f(x)|.
\]
由题设，对 $x_*$ 存在 $x'\in\lbrack a,b\rbrack$，使
\[
    m\leq |f(x')|\leq\dfrac{m}{2}.
\]
故 $m=0$，于是 $f(x_*)=0$。取 $x_0=x_*$ 即得结论。
\end{proof}

\begin{example}
设 $f\in C(\lbrack a,b\rbrack)$,若 $(a,b)$ 中无 $f(x)$ 的极值点,则 $f(x)$ 在 $\lbrack a,b\rbrack$ 上单调.
\end{example}

\begin{proof}[解]
反设 $f$ 在 $\lbrack a,b\rbrack$ 上不单调。连续函数在区间上若不是单调函数，则存在 $x_1<x_2<x_3$，使 $f(x_2)$ 不介于 $f(x_1)$ 与 $f(x_3)$ 之间。若
\[
    f(x_2)>\max\{f(x_1),f(x_3)\},
\]
则 $f$ 在 $\lbrack x_1,x_3\rbrack$ 上的最大值不可能在端点取得，故在 $(x_1,x_3)$ 内有极大值点；另一种不等式同理给出极小值点。两者都与 $(a,b)$ 中无极值点矛盾。因此 $f$ 必单调。
\end{proof}

\begin{example}
设 $f\in C(( -\infty,+\infty))$,若对任意的开区间 $(a,b)$,其值域 $f((a,b))$ 必是开区间,则 $f(x)$ 必单调.
\end{example}

\begin{proof}[解]
反设 $f$ 不单调。由连续性，存在 $x_1<x_2<x_3$，使 $f(x_2)$ 严格大于两端函数值，或严格小于两端函数值。不妨设为前一种。取包含 $x_2$ 而端点仍位于 $x_1,x_3$ 之间的开区间 $I$。$f$ 在其闭包上取得最大值，且该最大值在 $I$ 的内部取得。因此 $f(I)$ 含有自己的最大值，不可能是开区间，矛盾。故 $f$ 必单调；又常值函数把开区间映为单点，也不满足题设，所以实际上 $f$ 严格单调。
\end{proof}

\subsection{介值性、零点存在性}

\begin{example}
\begin{enumerate}[label=(\arabic*)]
    \item 设 $f(x)$ 在 $\lbrack a,b\rbrack$ 上单调,若 $f(x)$ 在 $\lbrack a,b\rbrack$ 上具有介值性,则 $f\in C(\lbrack a,b\rbrack)$;
    \item 设 $f(x)$ 定义在 $\lbrack a,b\rbrack$ 上,若对任意 $x\in(a,b)$,均存在极限 $\displaystyle\lim\limits_{t\to x}f(t)$,且又具有介值性,则 $f\in C(\lbrack a,b\rbrack)$.
\end{enumerate}
\end{example}

\begin{proof}[解]
第 (1) 问中，单调函数在一点不连续时，其左右极限之间存在一个非空跳跃区间，而该区间中的值不能由函数取得，这与介值性矛盾。因此 $f$ 在每个内点连续；端点处同理用单侧极限与介值性即可证明连续。

第 (2) 问中，固定 $x_0\in(a,b)$，记 $L=\displaystyle\lim\limits_{t\to x_0}f(t)$。若 $f(x_0)\ne L$，取严格介于二者之间的数 $c$。当 $t$ 充分接近但不等于 $x_0$ 时，$f(t)$ 全部位于 $L$ 的一个小邻域内，因而取不到 $c$；但介值性应用于 $x_0$ 与这样的 $t$，又要求函数在两点之间取得 $c$，矛盾。故 $f(x_0)=L$。端点处使用单侧版本作同样论证，遂有 $f\in C(\lbrack a,b\rbrack)$。
\end{proof}

\begin{example}
设 $f\in C(\lbrack a,b\rbrack)$,$x_i\in\lbrack a,b\rbrack\ (i=1,2,\ldots,n)$,则存在 $\xi\in\lbrack a,b\rbrack$,使得
\[
    f(\xi)=\dfrac{1}{n}\sum\limits_{i=1}^n f(x_i).
\]
\end{example}

\begin{proof}[解]
令
\[
    m=\min\limits_{1\leq i\leq n}f(x_i),
    \qquad
    M=\max\limits_{1\leq i\leq n}f(x_i).
\]
由每个 $f(x_i)$ 都位于 $m$ 与 $M$ 之间，对这些不等式求和再除以 $n$，得
\[
    m\leq\dfrac{1}{n}\sum\limits_{i=1}^n f(x_i)\leq M.
\]
取 $x_p,x_q$ 使 $f(x_p)=m$、$f(x_q)=M$。由连续函数的介值定理，在 $x_p$ 与 $x_q$ 之间存在 $\xi$，使
\[
    f(\xi)=\dfrac{1}{n}\sum\limits_{i=1}^n f(x_i).
\]
\end{proof}

\begin{example}
\begin{enumerate}[label=(\arabic*)]
    \item 设 $f\in C(\lbrack 0,1\rbrack)$,若 $f(0)=f(1)$,则存在 $\alpha,\beta$ 满足
    \[
        0\leq\alpha<\beta\leq1,
        \qquad \beta-\alpha=\dfrac{1}{5},
    \]
    使得 $f(\alpha)=f(\beta)$;
    \item 设 $f(x)$ 在 $(-\infty,+\infty)$ 上以 $T$ 为周期的连续函数,则存在 $\xi\in(-\infty,+\infty)$ 使得
    \[
        f\left(\xi+\dfrac{T}{2}\right)=f(\xi).
    \]
\end{enumerate}
\end{example}

\begin{proof}[解]
第 (1) 问令
\[
    g(x)=f\left(x+\dfrac{1}{5}\right)-f(x),
    \qquad 0\leq x\leq\dfrac{4}{5}.
\]
若某个 $g(k/5)=0$，结论立即成立；否则由
\[
    \sum\limits_{k=0}^{4}g\left(\dfrac{k}{5}\right)=f(1)-f(0)=0
\]
可知这些值中有正有负。由 $g$ 的连续性，存在 $\alpha\in\lbrack0,4/5\rbrack$ 使 $g(\alpha)=0$。令 $\beta=\alpha+1/5$ 即可。

第 (2) 问令
\[
    h(x)=f\left(x+\dfrac{T}{2}\right)-f(x).
\]
由周期性，$h(x+T/2)=-h(x)$。若 $h(x)=0$ 已得结论；否则 $h(x)$ 与 $h(x+T/2)$ 异号，由介值定理存在 $\xi$ 使 $h(\xi)=0$。
\end{proof}

\begin{example}
设 $f\in C(\lbrack 0,1\rbrack)$,$n\in\mathbb{N}$.
\begin{enumerate}[label=(\arabic*)]
    \item 若 $f(0)=f(1)$,则存在 $\xi_n\in\lbrack 0,\dfrac{1}{n}\rbrack$,使得
    \[
        f\left(\xi_n+\dfrac{1}{n}\right)=f(\xi_n);
    \]
    \item 若 $f(0)=0,f(1)=1$,则存在 $\xi_n\in(0,1)$,使得
    \[
        f\left(\xi_n+\dfrac{1}{n}\right)=f(\xi_n)+\dfrac{1}{n}.
    \]
\end{enumerate}
\end{example}

\begin{proof}[解]
第 (1) 问按当前题面不成立。取 $n=3$，并定义连续函数
\[
f(x)=
\begin{cases}
    3x, & 0\leq x\leq\dfrac{2}{3},\\
    6(1-x), & \dfrac{2}{3}\leq x\leq1.
\end{cases}
\]
则 $f(0)=f(1)=0$，但对每个 $\xi\in\lbrack0,1/3\rbrack$，均有
\[
    f\left(\xi+\dfrac{1}{3}\right)-f(\xi)=1,
\]
不存在题目所称的 $\xi_3$。

第 (2) 问需要先补充自然的定义域限制 $0\leq\xi\leq1-1/n$，且 $n\geq2$。在该修正下，令
\[
    h(x)=f\left(x+\dfrac{1}{n}\right)-f(x)-\dfrac{1}{n}.
\]
则
\[
    \sum\limits_{k=0}^{n-1}h\left(\dfrac{k}{n}\right)
    =f(1)-f(0)-1=0.
\]
因此这些网格点上或者已有 $h=0$，或者 $h$ 有正有负；由介值定理存在合法的 $\xi$ 使 $h(\xi)=0$。原题第 (1) 问及第 (2) 问的定义域写法均需核对原始资料。
\end{proof}

\begin{example}
设 $f\in C(( -\infty,+\infty))$,且以 $1$ 为周期,则存在 $\xi$,使得
\[
    f(\xi+\pi)=f(\xi).
\]
\end{example}

\begin{proof}[解]
令
\[
    g(x)=f(x+\pi)-f(x).
\]
取 $f$ 在一个周期 $\lbrack0,1\rbrack$ 上的最大值点 $x_0$。由周期性，$x_0$ 也可看作实轴上的最大值点，故
\[
    g(x_0)=f(x_0+\pi)-f(x_0)\leq0.
\]
另一方面，
\[
    g(x_0-\pi)=f(x_0)-f(x_0-\pi)\geq0.
\]
$g$ 连续，由介值定理，存在 $\xi\in\lbrack x_0-\pi,x_0\rbrack$ 使 $g(\xi)=0$，即 $f(\xi+\pi)=f(\xi)$。
\end{proof}

\begin{example}
设 $f,g\in C(\lbrack 0,1\rbrack)$,且
\[
    \max\limits_{x\in\lbrack 0,1\rbrack}f(x)=\max\limits_{x\in\lbrack 0,1\rbrack}g(x),
\]
试证:存在 $x_0\in\lbrack 0,1\rbrack$ 使得
\[
    e^{f(x_0)}+3f(x_0)=e^{g(x_0)}+3g(x_0).
\]
\end{example}

\begin{proof}[解]
记共同的最大值为 $M$，并取 $x_f,x_g\in\lbrack0,1\rbrack$，使 $f(x_f)=M$、$g(x_g)=M$。令 $h(x)=f(x)-g(x)$，则
\[
    h(x_f)=M-g(x_f)\geq0,
    \qquad
    h(x_g)=f(x_g)-M\leq0.
\]
由介值定理，存在 $x_0$ 使 $f(x_0)=g(x_0)$。函数 $\varphi(t)=e^t+3t$ 严格递增，故在此点
\[
    e^{f(x_0)}+3f(x_0)=e^{g(x_0)}+3g(x_0).
\]
\end{proof}

\begin{example}
设函数 $f(x)$ 在 $(-\infty,+\infty)$ 上连续,$n$ 为奇数,若
\[
    \lim\limits_{x\to+\infty}\dfrac{f(x)}{x^n}
    =\lim\limits_{x\to-\infty}\dfrac{f(x)}{x^n}=1,
\]
试证:方程 $f(x)+x^n=0$ 有实根.
\end{example}

\begin{proof}[解]
令 $F(x)=f(x)+x^n$。当 $x\to+\infty$ 时，
\[
    \dfrac{F(x)}{x^n}=\dfrac{f(x)}{x^n}+1\longrightarrow2,
\]
故 $F(x)\to+\infty$。当 $x\to-\infty$ 时同一商仍趋于 $2$，但因 $n$ 为奇数，$x^n\to-\infty$，故 $F(x)\to-\infty$。于是可取 $A<B$ 使 $F(A)<0<F(B)$。由连续性，存在 $\xi\in(A,B)$ 使 $F(\xi)=0$，即 $f(\xi)+\xi^n=0$。
\end{proof}

\begin{example}
设
\[
    f_n(x)=x+x^2+\cdots+x^n,
    \qquad n=2,3,\ldots.
\]
\begin{enumerate}[label=(\arabic*)]
    \item 证明:方程 $f_n(x)=1$ 在 $\lbrack 0,+\infty)$ 上有唯一的实根 $x_n$;
    \item 证明数列 $\{x_n\}$ 有极限,并求出 $\displaystyle\lim\limits_{n\to+\infty}x_n$.
\end{enumerate}
\end{example}

\begin{proof}[解]
在 $\lbrack0,+\infty)$ 上，$f_n$ 连续且严格递增，并且 $f_n(0)=0<1$、$f_n(1)=n>1$，故方程有唯一正根 $x_n\in(0,1)$。

又因 $f_{n+1}(x)>f_n(x)$（$x>0$），故 $x_{n+1}<x_n$。所以 $x_n$ 单调递减且有下界。对 $0<x<1$，
\[
    f_n(x)<\sum\limits_{k=1}^{\infty}x^k=\dfrac{x}{1-x}.
\]
在 $x=x_n$ 处左端为 $1$，故 $x_n>1/2$。设 $x_n\to L$，则 $1/2\leq L<1$，且 $x_n^n\to0$。由
\[
    1=\dfrac{x_n(1-x_n^n)}{1-x_n}
\]
取极限，得 $1=L/(1-L)$，所以 $L=1/2$。
\end{proof}

\begin{example}
设
\[
    f_n(x)=\cos x+\cos^2x+\cdots+\cos^n x,
\]
求证:
\begin{enumerate}[label=(\arabic*)]
    \item 对任意 $n$,$f_n(x)=1$ 在 $\left[0,\dfrac{\pi}{3}\right)$ 内有且仅有一根;
    \item 若 $x_n\in\left[0,\dfrac{\pi}{3}\right)$ 是 $f_n(x)=1$ 的根,则
    \[
        \lim\limits_{n\to\infty}x_n=\dfrac{\pi}{3}.
    \]
\end{enumerate}
\end{example}

\begin{proof}[解]
令 $t=\cos x$。当 $x\in\lbrack0,\pi/3)$ 时，$t\in(1/2,1\rbrack$，且
\[
    f_n(x)=t+t^2+\cdots+t^n
\]
随 $x$ 严格递减。在 $x=0$ 时其值为 $n>1$，而当 $x\to(\pi/3)^-$ 时其极限为
\[
    \sum\limits_{k=1}^n\dfrac{1}{2^k}=1-\dfrac{1}{2^n}<1.
\]
故存在唯一根 $x_n\in\lbrack0,\pi/3)$。

记 $t_n=\cos x_n$。由 $\displaystyle\sum\limits_{k=1}^n t_n^k=1$ 可知 $t_n>1/2$；随着 $n$ 增大，根 $t_n$ 严格递减，故 $t_n\to L\geq1/2$。又 $t_n\leq t_2<1$，所以 $t_n^n\to0$。在
\[
    1=\dfrac{t_n(1-t_n^n)}{1-t_n}
\]
中取极限，得 $L=1/2$。因此
\[
    x_n=\arccos t_n\longrightarrow\arccos\dfrac{1}{2}=\dfrac{\pi}{3}.
\]
\end{proof}

\begin{example}
设 $f(x)$ 在 $(0,+\infty)$ 连续有界,试证:对任意 $T$,存在 $x_n\to+\infty$,使得
\[
    \lim\limits_{n\to\infty}[f(x_n+T)-f(x_n)]=0.
\]
\end{example}

\begin{proof}[解]
只需考虑 $T>0$；当 $T=0$ 时差值恒为 $0$；当 $T<0$ 时，可先对 $-T>0$ 取到相应点列，再将点列平移 $-T$ 即可。下设 $T>0$。反设结论不成立，则存在 $\varepsilon_0>0$ 与 $R>0$，使所有 $x\geq R$ 都满足
\[
    |f(x+T)-f(x)|\geq\varepsilon_0.
\]
连续函数 $f(x+T)-f(x)$ 在连通区间 $\lbrack R,+\infty)$ 上不能变号，否则必有零点。因此它恒不小于 $\varepsilon_0$，或恒不大于 $-\varepsilon_0$。前一种情形下
\[
    f(R+nT)-f(R)
    =\sum\limits_{k=0}^{n-1}[f(R+(k+1)T)-f(R+kT)]
    \geq n\varepsilon_0,
\]
与 $f$ 有界矛盾；后一种情形同理。故对每个正整数 $n$，均可取 $x_n\geq n$ 使
\[
    |f(x_n+T)-f(x_n)|<\dfrac{1}{n},
\]
结论成立。
\end{proof}

\begin{example}
设 $f$ 在 $\lbrack 0,n\rbrack$ 上连续,$f(0)=f(n)$,试证:至少存在 $n$ 组不同的解 $(x,y)$ 使得 $f(x)=f(y)$,且 $y-x>0$ 为整数.
\end{example}

\begin{proof}[解]
对 $t\in\lbrack0,1\rbrack$，考虑 $n$ 条连续曲线
\[
    F_k(t)=f(t+k),
    \qquad k=0,1,\ldots,n-1.
\]
在 $t=0$ 时，它们的端点值依次为 $f(0),f(1),\ldots,f(n-1)$；在 $t=1$ 时，端点值依次变为 $f(1),\ldots,f(n-1),f(n)=f(0)$，即作了一次循环置换。

一组连续曲线的上下次序若发生一次循环置换，至少需要 $n-1$ 次两两相交：把每次相交看作允许相邻次序交换，则一个 $n$-循环至少由 $n-1$ 个换位组成；若多条曲线在同一点相交，该点同时给出相应的多个不同曲线对，计数只会增加。因此至少有 $n-1$ 个不同的曲线对 $(i,j)$ 与相交参数 $t$，满足
\[
    F_i(t)=F_j(t).
\]
这给出 $f(t+i)=f(t+j)$，且两点距离 $j-i$ 为正整数。再加上题设给出的不同解对 $(0,n)$，总计至少有 $n$ 组不同的解。
\end{proof}

\subsection{一致连续性}

\begin{example}
判断下列函数在指定区间上的一致连续性:
\begin{enumerate}[label=(\arabic*)]
    \item $f(x)=\sin^2x$,$\lbrack 0,+\infty)$;
    \item $f(x)=\sin x^2$,$\lbrack 0,+\infty)$;
    \item $f(x)=\sqrt{x}$,$\lbrack 0,+\infty)$;
    \item $f(x)=\dfrac{1}{x}$,$(0,1)$;
    \item $g(x)=x^2$,$(1,+\infty)$.
\end{enumerate}
\end{example}

\begin{proof}[解]
第 (1) 问中，$|f'(x)|=|\sin2x|\leq1$，故 $f$ 满足 Lipschitz 条件，从而一致连续。

第 (2) 问不一致连续。取
\[
    x_n=\sqrt{2n\pi+\dfrac{\pi}{2}},
    \qquad
    y_n=\sqrt{2n\pi+\dfrac{3\pi}{2}}.
\]
则 $|x_n-y_n|\to0$，但 $|\sin x_n^2-\sin y_n^2|=2$。

第 (3) 问一致连续，因为
\[
    |\sqrt{x}-\sqrt{y}|\leq\sqrt{|x-y|}.
\]

第 (4) 问不一致连续。取 $x_n=1/n$、$y_n=1/(n+1)$，则 $|x_n-y_n|\to0$，但 $|1/x_n-1/y_n|=1$。

第 (5) 问不一致连续。取 $x_n=n$、$y_n=n+1/n$，则 $|x_n-y_n|\to0$，而
\[
    |x_n^2-y_n^2|=2+\dfrac{1}{n^2}\not\to0.
\]
\end{proof}

\begin{example}
判断下列函数在指定区间上的一致连续性:
\begin{enumerate}[label=(\arabic*)]
    \item $f(x)=\ln x$,$(0,1\rbrack$;
    \item $f(x)=\dfrac{\sin x}{x}$,$(0,1\rbrack$;
    \item $f(x)=e^x\cos\dfrac{1}{x}$,$(0,1\rbrack$;
    \item $f(x)=\cos x\cdot\cos\dfrac{\pi}{x}$,$(0,1)$;
    \item $f(x)=x\sin\dfrac{1}{x}$,$(0,+\infty)$;
    \item $f(x)=\dfrac{x+2}{x+1}\sin\dfrac{1}{x}$,$(0,1)$ 及 $(1,+\infty)$.
\end{enumerate}
\end{example}

\begin{proof}[解]
第 (1) 问不一致连续，因为 $x\to0^+$ 时 $\ln x\to-\infty$，而一致连续函数在有限端点附近必有有限单侧极限。

第 (2) 问中，令 $f(0)=1$，则 $\sin x/x$ 连续延拓到紧区间 $\lbrack0,1\rbrack$，故一致连续。

第 (3) 问不一致连续。取 $x_n=1/(2n\pi)$、$y_n=1/((2n+1)\pi)$，则 $|x_n-y_n|\to0$，而函数值分别趋于 $1$ 与 $-1$。

第 (4) 问不一致连续。取 $x_n=1/(2n)$、$y_n=1/(2n+1)$，则两点距离趋于零，而 $\cos(\pi/x_n)=1$、$\cos(\pi/y_n)=-1$，且 $\cos x_n,\cos y_n\to1$。

第 (5) 问一致连续。在 $x=0$ 处补定义为 $0$ 后，函数在 $\lbrack0,1\rbrack$ 上连续；在 $\lbrack1,+\infty)$ 上其导数
\[
    \sin\dfrac{1}{x}-\dfrac{1}{x}\cos\dfrac{1}{x}
\]
有界，故两段可合并为一致连续。

第 (6) 问在 $(0,1)$ 上不一致连续，因为系数 $(x+2)/(x+1)\to2$，而 $\sin(1/x)$ 在零点附近保持全振荡；在 $(1,+\infty)$ 上导数有界，故一致连续。
\end{proof}

\begin{example}
判断下列函数 $f(x)$ 在指定区间上的一致连续性:
\begin{enumerate}[label=(\arabic*)]
    \item $f(x)=\dfrac{x^3}{x+1}\sin\dfrac{1}{x}$,$(0,+\infty)$;
    \item $f(x)=\sqrt{\dfrac{x^4+1}{x^2+1}}$,$(-\infty,+\infty)$;
    \item $f(x)=\sqrt[3]{x^3-x^2-x+1}$,$(-\infty,+\infty)$.
\end{enumerate}
\end{example}

\begin{proof}[解]
三问中的函数都一致连续。

第 (1) 问在 $x=0$ 处补定义为 $0$ 后连续；其导数在 $(0,+\infty)$ 上连续，且当 $x\to0^+$ 时函数为 $O(x^2)$，当 $x\to+\infty$ 时函数为 $x+O(1)$，相应导数在两端均有界，故全区间上满足某个 Lipschitz 条件。

第 (2) 问记
\[
    f(x)=\sqrt{\dfrac{x^4+1}{x^2+1}}.
\]
$f'$ 连续，且当 $x\to\pm\infty$ 时 $f'(x)\to\pm1$，故 $f'$ 在实轴上有界，因而 $f$ 一致连续。

第 (3) 问中
\[
    x^3-x^2-x+1=(x-1)^2(x+1).
\]
函数仅在 $x=-1,1$ 附近导数可能无界，但立方根满足
\[
    |\sqrt[3]{u}-\sqrt[3]{v}|\leq\sqrt[3]{4|u-v|},
\]
故在这两个有限邻域内仍一致连续；在其补集上导数连续有界，且在无穷远处导数趋于 $1$。有限拼接后，函数在实轴上一致连续。
\end{proof}

\begin{example}
设 $f$ 在 $\lbrack 0,+\infty)$ 上满足 Lipschitz 条件,证明:$f(x^\alpha)\ (0<\alpha<1)$ 在 $\lbrack 0,+\infty)$ 上一致连续.
\end{example}

\begin{proof}[解]
设 $f$ 的 Lipschitz 常数为 $L$。当 $x,y\geq0$ 且 $0<\alpha<1$ 时，幂函数满足
\[
    |x^\alpha-y^\alpha|\leq|x-y|^\alpha.
\]
因此
\[
    |f(x^\alpha)-f(y^\alpha)|
    \leq L|x^\alpha-y^\alpha|
    \leq L|x-y|^\alpha.
\]
若 $L=0$，则 $f$ 为常值函数，结论直接成立。若 $L>0$，给定 $\varepsilon>0$，取
$\delta=(\varepsilon/L)^{1/\alpha}$，即可得到一致连续性。
\end{proof}

\begin{example}[区间合并]
\begin{enumerate}[label=(\arabic*)]
    \item 若函数 $f$ 在区间 $(a,b\rbrack$ 和 $\lbrack b,c)$ 上分别一致连续,证明:$f$ 在 $(a,c)$ 上一致连续;
    \item 若函数 $f$ 在区间 $(a,b)$ 和 $\lbrack b,c)$ 上分别一致连续,此时 $f$ 在 $(a,c)$ 上是否一致连续?
    \item 证明:函数 $f(x)=\dfrac{|\sin x|}{x}$ 在区间 $(-1,0)$ 和 $(0,1)$ 均一致连续,但在 $(-1,0)\cup(0,1)$ 上非一致连续.
\end{enumerate}
\end{example}

\begin{proof}[解]
第 (1) 问成立。给定 $\varepsilon>0$，分别取两段的一致连续参数 $\delta_1,\delta_2$。令 $\delta=\min\{\delta_1,\delta_2\}$。若 $x,y$ 位于同一段，直接应用对应估计；若 $x<b<y$ 且 $|x-y|<\delta$，则
\[
    |f(x)-f(y)|
    \leq|f(x)-f(b)|+|f(b)-f(y)|<2\varepsilon.
\]
将开始时的误差改取为 $\varepsilon/2$ 即得结论。

第 (2) 问不一定成立。例如令 $f(x)=0$（$a<x<b$），而令 $f(x)=1$（$b\leq x<c$）。它在两段上分别一致连续，但在 $b$ 处不连续，故在 $(a,c)$ 上不一致连续。

第 (3) 问中，函数在 $x\to0^-$ 与 $x\to0^+$ 时分别趋于 $-1$ 与 $1$，所以它可分别连续延拓到两个紧区间，因而在 $(-1,0)$ 与 $(0,1)$ 上分别一致连续。但取 $x_n=-1/n$、$y_n=1/n$，则 $|x_n-y_n|\to0$，而函数值之差趋于 $2$，故在并集上不一致连续。
\end{proof}

\begin{example}
设 $f(x)$ 在 $\lbrack 0,+\infty)$ 上一致连续,且有
\[
    \lim\limits_{n\to\infty}f(x+n)=0,
    \qquad \forall x\in\lbrack 0,+\infty).
\]
则
\[
    \lim\limits_{x\to+\infty}f(x)=0.
\]
\end{example}

\begin{proof}[解]
对 $n\in\mathbb{N}$ 定义
\[
    g_n(t)=f(n+t),
    \qquad t\in\lbrack0,1\rbrack.
\]
由 $f$ 的一致连续性，函数族 $\{g_n\}$ 等度连续；题设又说明对每个固定 $t$，$g_n(t)\to0$。给定 $\varepsilon>0$，取 $\delta>0$，使 $|u-v|<\delta$ 时 $|f(u)-f(v)|<\varepsilon/2$。在 $\lbrack0,1\rbrack$ 中取有限个 $\delta$-网点 $t_1,\ldots,t_m$。对每个网点，存在 $N_j$ 使 $n\geq N_j$ 时 $|g_n(t_j)|<\varepsilon/2$。令 $N=\max\{N_1,\ldots,N_m\}$，则 $n\geq N$ 时，对任意 $t$ 选取满足 $|t-t_j|<\delta$ 的网点，有
\[
    |g_n(t)|\leq|g_n(t)-g_n(t_j)|+|g_n(t_j)|<\varepsilon.
\]
所以 $g_n\to0$ 在 $\lbrack0,1\rbrack$ 上一致。任意充分大的 $x$ 可写为 $x=n+t$，其中 $t\in\lbrack0,1)$，故 $f(x)\to0$。
\end{proof}

\begin{example}
设 $f(x)$ 在区间 $I$ 上有定义,记
\[
    w_f(\delta)
    =\sup\limits_{\substack{x',x''\in I\\|x'-x''|<\delta}}
    |f(x')-f(x'')|,
\]
则 $f$ 在 $I$ 上一致连续的充要条件是
\[
    \lim\limits_{\delta\to0^+}w_f(\delta)=0.
\]
\end{example}

\begin{proof}[解]
若 $f$ 在 $I$ 上一致连续，则对任意 $\varepsilon>0$，存在 $\delta_0>0$，使 $x',x''\in I$ 且 $|x'-x''|<\delta_0$ 时
\[
    |f(x')-f(x'')|<\varepsilon.
\]
因此对 $0<\delta\leq\delta_0$ 有 $w_f(\delta)\leq\varepsilon$，从而 $w_f(\delta)\to0$。

反之，若 $w_f(\delta)\to0$，给定 $\varepsilon>0$，可取 $\delta>0$ 使 $w_f(\delta)<\varepsilon$。于是任意 $x',x''\in I$ 只要 $|x'-x''|<\delta$，便有
\[
    |f(x')-f(x'')|\leq w_f(\delta)<\varepsilon.
\]
这正是 $f$ 在 $I$ 上一致连续的定义。
\end{proof}

\end{document}

